\documentclass[11pt,a4paper]{article}
\usepackage[utf8]{inputenc}
\usepackage{amsmath,amssymb,amsthm}
\usepackage{algorithm}
\usepackage{algpseudocode}
\usepackage{booktabs}
\usepackage{hyperref}
\usepackage{geometry}
\geometry{margin=1in}

\title{Recurrent Neural Networks (RNN) and Backpropagation Through Time (BPTT): Mathematical Formulations and Gradient Vanishing Dynamics}
\author{Team Scaibu \\ \small Email: \texttt{jaydeep.9664920749@gmail.com} \\ \small Architecture and Core Engine Team}
\date{\today}

\newtheorem{theorem}{Theorem}
\newtheorem{lemma}[theorem]{Lemma}
\newtheorem{proposition}[theorem]{Proposition}
\theoremstyle{definition}
\newtheorem{definition}{Definition}
\newtheorem{remark}{Remark}

\begin{document}

\maketitle

\begin{abstract}
Recurrent Neural Networks (RNNs) parameterize sequential temporal state transitions through stationary weight sharing across discrete time steps. Training recurrent models over extended temporal horizons requires unrolling the recurrence graph across $T$ steps, computing exact analytical gradients via Backpropagation Through Time (BPTT). This document formalizes the forward unrolling and backward adjoint state equations, mathematically proves the exponential vanishing and exploding gradient phenomenon via singular value spectrum analysis of the recurrent Jacobian, presents full pseudocode, and delivers a verified step-by-step numerical calculation.
\end{abstract}

\section{Notation and Mathematical Preliminaries}

Let $\mathbf{x} = (\mathbf{x}_1, \dots, \mathbf{x}_T)$ denote an input sequence where $\mathbf{x}_t \in \mathbb{R}^{d_x}$, $\mathbf{h}_t \in \mathbb{R}^{d_h}$ denote the hidden state vector at time $t$, and $\mathbf{y}_t \in \mathbb{R}^{d_y}$ denote the output prediction.

\begin{table}[h]
\centering
\caption{Mathematical Notation for RNN and BPTT}
\begin{tabular}{lll}
\toprule
\textbf{Symbol} & \textbf{Domain} & \textbf{Semantic Description} \\
\midrule
$W_h$ & $\mathbb{R}^{d_h \times d_h}$ & Recurrent hidden-to-hidden transition weight matrix \\
$W_x$ & $\mathbb{R}^{d_h \times d_x}$ & Input-to-hidden projection weight matrix \\
$W_y$ & $\mathbb{R}^{d_y \times d_h}$ & Hidden-to-output emission weight matrix \\
$\mathbf{b}_h, \mathbf{b}_y$ & $\mathbb{R}^{d_h}, \mathbb{R}^{d_y}$ & Hidden and emission bias vectors \\
$\mathbf{a}_t$ & $\mathbb{R}^{d_h}$ & Pre-activation vector at time $t$: $\mathbf{a}_t = W_h \mathbf{h}_{t-1} + W_x \mathbf{x}_t + \mathbf{b}_h$ \\
$\boldsymbol{\delta}_t$ & $\mathbb{R}^{d_h}$ & Adjoint backpropagation error vector: $\frac{\partial \mathcal{L}}{\partial \mathbf{a}_t}$ \\
$\mathcal{L}$ & $\mathbb{R}$ & Cumulative sequence loss: $\sum_{t=1}^T \mathcal{L}_t(\mathbf{y}_t, \mathbf{y}_t^*)$ \\
\bottomrule
\end{tabular}
\end{table}

\section{Problem Definition and Core Concepts}

\subsection{3-Level Explanation}
\begin{enumerate}
    \item \textbf{Intuitive (High School level):} An RNN is like a person reading a book one word at a time, keeping a running memory in their head. To train the reader, we replay the entire chapter backwards in time to see which word caused an error, adjusting the reader's memory habits.
    \item \textbf{Engineering Level:} An RNN shares the same weight matrices $(W_h, W_x, W_y)$ across all $T$ time steps. BPTT unrolls the computation into a $T$-layer feedforward network with shared parameters, propagating gradients backward along the temporal chain and accumulating parameter updates across all steps.
    \item \textbf{Mathematical Level:} Given sequence $\{\mathbf{x}_t\}_{t=1}^T$, forward recurrence satisfies:
    \begin{equation}
    \mathbf{h}_t = \tanh(W_h \mathbf{h}_{t-1} + W_x \mathbf{x}_t + \mathbf{b}_h), \quad \mathbf{y}_t = W_y \mathbf{h}_t + \mathbf{b}_y
    \end{equation}
    and backward adjoint recurrence is given by:
    \begin{equation}
    \frac{\partial \mathcal{L}}{\partial \mathbf{h}_t} = W_y^T \frac{\partial \mathcal{L}_t}{\partial \mathbf{y}_t} + W_h^T \boldsymbol{\delta}_{t+1}, \quad \boldsymbol{\delta}_t = \frac{\partial \mathcal{L}}{\partial \mathbf{h}_t} \odot (1 - \mathbf{h}_t^2)
    \end{equation}
\end{enumerate}

\subsection{5-Step Algorithmic Framework}
\begin{enumerate}
    \item \textbf{Temporal Forward Unrolling:} Compute sequence of hidden states $\mathbf{h}_1, \dots, \mathbf{h}_T$ and outputs $\mathbf{y}_1, \dots, \mathbf{y}_T$.
    \item \textbf{Loss Evaluation:} Compute per-step loss $\mathcal{L}_t$ and terminal output gradient $\frac{\partial \mathcal{L}_T}{\partial \mathbf{y}_T}$.
    \item \textbf{Backward Adjoint Propagation:} Iterate backwards $t = T, \dots, 1$, evaluating adjoint vectors $\boldsymbol{\delta}_t$.
    \item \textbf{Parameter Gradient Accumulation:} Sum outer product gradients $\frac{\partial \mathcal{L}}{\partial W_h} = \sum_t \boldsymbol{\delta}_t \mathbf{h}_{t-1}^T$, $\frac{\partial \mathcal{L}}{\partial W_x} = \sum_t \boldsymbol{\delta}_t \mathbf{x}_t^T$.
    \item \textbf{Optimizer Parameter Update:} Apply gradient descent updates to shared parameters.
\end{enumerate}

\section{Algorithm Specification}

\begin{algorithm}[H]
\caption{Backpropagation Through Time (BPTT) for Vanilla RNN}
\label{alg:bptt}
\begin{algorithmic}[1]
\Require Sequence $\{\mathbf{x}_t\}_{t=1}^T$, targets $\{\mathbf{y}_t^*\}_{t=1}^T$, weights $W_x, W_h, W_y$, biases $\mathbf{b}_h, \mathbf{b}_y$.
\Ensure Gradients $\nabla_{W_x} \mathcal{L}, \nabla_{W_h} \mathcal{L}, \nabla_{W_y} \mathcal{L}, \nabla_{\mathbf{b}_h} \mathcal{L}, \nabla_{\mathbf{b}_y} \mathcal{L}$.
\State $\mathbf{h}_0 \gets \mathbf{0}$
\For{$t = 1$ \textbf{to} $T$} \Comment{Forward Pass}
    \State $\mathbf{a}_t \gets W_h \mathbf{h}_{t-1} + W_x \mathbf{x}_t + \mathbf{b}_h$
    \State $\mathbf{h}_t \gets \tanh(\mathbf{a}_t)$
    \State $\mathbf{y}_t \gets W_y \mathbf{h}_t + \mathbf{b}_y$
\EndFor
\State Initialize gradient accumulators $\nabla W_x \gets \mathbf{0}, \nabla W_h \gets \mathbf{0}, \nabla W_y \gets \mathbf{0}, \nabla \mathbf{b}_h \gets \mathbf{0}, \nabla \mathbf{b}_y \gets \mathbf{0}$
\State $\mathbf{d}\mathbf{h}_{\text{next}} \gets \mathbf{0}$
\For{$t = T$ \textbf{downto} $1$} \Comment{Backward Pass}
    \State $\mathbf{d}\mathbf{y}_t \gets \mathbf{y}_t - \mathbf{y}_t^*$
    \State $\nabla W_y \gets \nabla W_y + \mathbf{d}\mathbf{y}_t \mathbf{h}_t^T, \quad \nabla \mathbf{b}_y \gets \nabla \mathbf{b}_y + \mathbf{d}\mathbf{y}_t$
    \State $\mathbf{d}\mathbf{h}_t \gets W_y^T \mathbf{d}\mathbf{y}_t + \mathbf{d}\mathbf{h}_{\text{next}}$
    \State $\boldsymbol{\delta}_t \gets \mathbf{d}\mathbf{h}_t \odot (1 - \mathbf{h}_t^2)$
    \State $\nabla W_h \gets \nabla W_h + \boldsymbol{\delta}_t \mathbf{h}_{t-1}^T$
    \State $\nabla W_x \gets \nabla W_x + \boldsymbol{\delta}_t \mathbf{x}_t^T$
    \State $\nabla \mathbf{b}_h \gets \nabla \mathbf{b}_h + \boldsymbol{\delta}_t$
    \State $\mathbf{d}\mathbf{h}_{\text{next}} \gets W_h^T \boldsymbol{\delta}_t$
\EndFor
\State \Return Gradients
\end{algorithmic}
\end{algorithm}

\section{Theoretical Guarantees and Gradient Vanishing Dynamics}

\begin{theorem}[Exponential Vanishing and Exploding Gradient Theorem]
\label{thm:vanishing_gradient}
Let the temporal gradient Jacobian be defined by $\frac{\partial \mathbf{h}_T}{\partial \mathbf{h}_t} = \prod_{k=t+1}^T \frac{\partial \mathbf{h}_k}{\partial \mathbf{h}_{k-1}} = \prod_{k=t+1}^T \text{diag}(1 - \mathbf{h}_k^2) W_h$.
Let $\sigma_{\max}(W_h) = \lambda_1$.
\begin{enumerate}
    \item If $\lambda_1 < 1$, then $\left\| \frac{\partial \mathbf{h}_T}{\partial \mathbf{h}_t} \right\|_2 \le \lambda_1^{T-t} \xrightarrow{T-t \to \infty} 0$ (Exponential Vanishing).
    \item If $\lambda_1 > 1$, the gradient norm can grow as $\Omega(\lambda_1^{T-t})$ along principal singular vectors (Exponential Exploding).
\end{enumerate}
\end{theorem}

\begin{proof}
By sub-multiplicativity of matrix operator 2-norms:
\begin{equation}
\left\| \frac{\partial \mathbf{h}_T}{\partial \mathbf{h}_t} \right\|_2 \le \prod_{k=t+1}^T \|\text{diag}(1 - \mathbf{h}_k^2)\|_2 \|W_h\|_2
\end{equation}
Because $|1 - h_{k, i}^2| \le 1$ for all $h \in (-1, 1)$, $\|\text{diag}(1 - \mathbf{h}_k^2)\|_2 \le 1$. Hence:
\begin{equation}
\left\| \frac{\partial \mathbf{h}_T}{\partial \mathbf{h}_t} \right\|_2 \le \|W_h\|_2^{T-t} = \lambda_1^{T-t}
\end{equation}
For $\lambda_1 < 1$, this upper bound contracts exponentially to 0 as temporal distance $(T - t) \to \infty$.
\end{proof}

\section{Complexity Analysis}

\begin{itemize}
    \item \textbf{Time Complexity:} $\mathcal{O}(T \cdot (d_h^2 + d_h d_x + d_y d_h))$ FLOPs per sequence (strictly sequential across time $T$).
    \item \textbf{Space Complexity:} $\mathcal{O}(T \cdot d_h)$ memory to cache the hidden state trajectory $\{\mathbf{h}_t\}_{t=0}^T$ for backward pass evaluation.
\end{itemize}

\section{Exactness and Error Analysis}

BPTT provides the exact mathematical gradient of the cumulative unrolled loss function without any approximation error. However, due to finite 32-bit floating point precision, vanishing gradients underflow to exact zero for $(T - t) > 50$.

\section{Worked Numerical Example}

Let $d_x=1, d_h=1, d_y=1, T=2$.
Parameters: $W_x = [[1.0]], W_h = [[0.5]], b_h = [0.0], W_y = [[1.0]], b_y = [0.0]$.
Inputs: $x_1 = 1.0, x_2 = 1.0$; Targets: $y_1^* = 0.0, y_2^* = 0.0$.
Initial $h_0 = 0.0$.

\textbf{Forward Pass:}
\begin{itemize}
    \item $t=1$: $a_1 = 0.5(0.0) + 1.0(1.0) = 1.0 \implies h_1 = \tanh(1.0) = 0.7616$. $y_1 = 1.0(0.7616) = 0.7616$.
    \item $t=2$: $a_2 = 0.5(0.7616) + 1.0(1.0) = 1.3808 \implies h_2 = \tanh(1.3808) = 0.8807$. $y_2 = 0.8807$.
\end{itemize}

\textbf{Backward Pass:}
\begin{itemize}
    \item $t=2$: $dy_2 = 0.8807 - 0 = 0.8807$.
    $dh_2 = 1.0(0.8807) = 0.8807$.
    $\delta_2 = 0.8807 \times (1 - 0.8807^2) = 0.8807 \times (1 - 0.7756) = 0.8807 \times 0.2244 = 0.1976$.
    $dh_{\text{next}} = 0.5 \times 0.1976 = 0.0988$.
    \item $t=1$: $dy_1 = 0.7616 - 0 = 0.7616$.
    $dh_1 = 1.0(0.7616) + 0.0988 = 0.8604$.
    $\delta_1 = 0.8604 \times (1 - 0.7616^2) = 0.8604 \times 0.4200 = 0.3614$.
\end{itemize}

Accumulated $\frac{\partial \mathcal{L}}{\partial W_h} = \delta_2 h_1 + \delta_1 h_0 = 0.1976(0.7616) + 0 = 0.1505$.

\section{Numerical Considerations and Edge Cases}

\begin{itemize}
    \item \textbf{Gradient Clipping by Norm:} Enforcing $\|\mathbf{g}\| \gets \min(1, \frac{\tau}{\|\mathbf{g}\|}) \mathbf{g}$ is strictly necessary to prevent exploding gradients when spectral radius $\rho(W_h) > 1$.
    \item \textbf{Unitary / Orthogonal Initialization:} Initializing $W_h$ as an orthogonal matrix ($W_h^T W_h = I$) preserves gradient norms during early training.
\end{itemize}

\section{References}
\begin{enumerate}
    \item Werbos, P. J. (1990). Backpropagation through time: what it does and how to do it. \emph{Proceedings of the IEEE}, 78(10), 1550--1560.
    \item Pascanu, R., Mikolov, T., & Bengio, Y. (2013). On the difficulty of training recurrent neural networks. \emph{ICML}, 1310--1318.
\end{enumerate}

\end{document}
